% Current mathematical front matter.  The historical review pages that follow
% are preserved as an audit record and are not silently rewritten.
\pdfbookmark[0]{Current review findings and proof status}{audit-corrective-addendum}
\thispagestyle{empty}
\begingroup
\raggedright
\setlength{\parindent}{0pt}
\setlength{\parskip}{0pt}
\setlength{\fboxsep}{5pt}
\newcommand{\ResultHead}[1]{\par\vspace{3.4pt}{\fontsize{9.15pt}{10.25pt}\selectfont\bfseries\color{AuditReviewed}#1}\par\vspace{1.2pt}}
\newcommand{\ResultSmall}{\fontsize{8.05pt}{9.23pt}\selectfont}

{\fontsize{15.8pt}{17.2pt}\selectfont\bfseries Independent Pro-mode reassessment: the mathematical findings}\par
\vspace{1pt}
{\fontsize{7.75pt}{8.9pt}\selectfont\color{AuditProvenance}A cold-entry account of Reviewer B's corrective audit, followed on page 3 by results proved subsequently in the Codex validate/repair/disprove workbench.}
\vspace{4pt}

\noindent\fcolorbox{AuditExact}{AuditPale}{%
\begin{minipage}{0.958\textwidth}
\ResultSmall
\textbf{Reviewer B's exact revised disposition: RECONSIDERED --- PROVISIONALLY
FAVORABLE WITH FOCUSED VERIFICATION REQUESTS.}
The earlier topology-based rejection is withdrawn.  Reviewer B found no
load-bearing defect in the intrinsic nearby-cycle/Leray and based van Kampen
route; it did not claim completed journal certification.
\end{minipage}}

\ResultHead{1. The correction is a distinction between two morphisms}
{\ResultSmall
Let $F$ be a nearby smooth torus fibre and $W$ the reduced cusp fibre.
The optional explicit collapse $c:F\to W$ would induce
\[
 c^*:H^q(W;\mathbb Z)\longrightarrow H^q(F;\mathbb Z),
\]
whereas nearby cycles supply the intrinsic stalk specialization
\[
 \operatorname{sp}^q:H^q(W;\mathbb Z)
 \longrightarrow H^q(F;\mathbb Z)^{T_0}.
\]
The earlier audit treated the uncompleted identification of $c^*$ with
specialization as a premise of both topology proofs.  Section~7.7 uses only
\(\operatorname{sp}^q\).  Reviewer B reconstructed Appendix~B's proof,
independent of $c$, that
\[
 \operatorname{sp}^q:H^q(W;\mathbb Z)
 \xrightarrow{\ \sim\ }H^q(F;\mathbb Z)^{T_0},\qquad 0\le q\le4,
 \qquad \operatorname{rk}=(1,2,4,2,1).
\]
Thus incomplete certification of the collapse map does not block the second
topology route.
}

\ResultHead{2. The same integer controls Leray and the based fundamental group}
{\ResultSmall
For the manuscript's twists
\[
 (\ell_0,\ell_1,\ell_2)=(0,1,-1),\qquad
 p=12\ell_0-4\ell_1-3\ell_2=-1.
\]
The reconstructed global stalk sections are
\[
 H^0(R^1f_*\mathbb Z)=\mathbb Z(12\gamma),\quad
 H^0(R^2f_*\mathbb Z)=\mathbb Z(2q),\quad
 H^0(R^3f_*\mathbb Z)=\mathbb Z(2\gamma uw),
\]
and Reviewer B obtained, up to a common orientation sign,
\[
 d_2(12\gamma)=\pm p\omega,\qquad
 d_2(2q)=\pm p\,\omega\bar\delta,\qquad
 d_2(2\gamma uw)=\pm p\,\omega[\gamma\delta].
\]
Because $p=-1$, all three rank-one maps are units, giving
\(H^1(X;\mathbb Z)=H^2(X;\mathbb Z)=H^3(X;\mathbb Z)=0\).  The based
overlap calculation gives
\[
 \pi_1(X)=\langle c,x,y\mid c\ {\rm central},\
 x^3=c^{\ell_1},\ y^4=c^{\ell_2},\ xy=c^{\ell_0}\rangle;
\]
its nonabelian elimination is cyclic of order $|p|$, hence trivial for
the selected twists.  Recognition then gives the integral homology and
smooth type of $S^6$.
}

\ResultHead{3. The algebraic-dimension computation is independent of topology}
{\ResultSmall
Reviewer B also reconstructed the period argument.  Off a countable set the
integral $(1,1)$-kernel is
\[
 \operatorname{NS}(F_b)=\mathbb Z\eta,\qquad
 \eta=uw+6\gamma\delta,\qquad \eta^2=12,
\]
and the corresponding Hermitian form has
\[
 \det G_\eta=\frac23\operatorname{Im}\tau
 \left(\operatorname{Im}\beta-
 \frac{6(\operatorname{Im}\mu)^2}{\operatorname{Im}\tau}\right)<0.
\]
The torus algebraic-reduction alternatives are thereby excluded; meromorphic
descent yields
\[
 \mathcal M(X)=f^*\mathbb C(t),\qquad a(X)=1,
\]
and identifies $f:X\to\mathbb P^1$ as the algebraic reduction, without
using Section~7 or the CDP comparison.
}

\vfill
{\fontsize{6.7pt}{7.65pt}\selectfont\color{AuditProvenance}
\textbf{Scope of the independent review.} Reviewer B returned to the
canonical source and independently examined the reconstruction evidence
existing up to its evidentiary cutoff.  That cutoff preceded approximately
one and a half hours of later Codex work.  Its audit is therefore neither
obsolete nor evidence for arguments it did not see: it remains an
independent comparison against the subsequent proofs on page~3.}
\endgroup

\newpage
\thispagestyle{empty}
\begingroup
\raggedright
\setlength{\parindent}{0pt}
\setlength{\parskip}{0pt}
\setlength{\fboxsep}{5pt}
\newcommand{\ResultHead}[1]{\par\vspace{3.7pt}{\fontsize{9.2pt}{10.3pt}\selectfont\bfseries\color{AuditReviewed}#1}\par\vspace{1.3pt}}
\newcommand{\ResultSmall}{\fontsize{8.15pt}{9.42pt}\selectfont}

{\fontsize{15.2pt}{16.6pt}\selectfont\bfseries What Reviewer B actually checked}
\par\vspace{2pt}
{\ResultSmall
This page states the review's proof-level scope, not merely its conclusion.
The accompanying evidence contains a revised theorem ledger, a revised risk
ledger, a 17/17 integration harness, and the earlier 38/38 exact core harness.
Those counts are receipts for named suites, not 55 independent theorems.
}

\ResultHead{Appendix B: integral nearby-cycle specialization}
{\ResultSmall
The reconstruction checks the local models
\(t=z_0\), \(t=z_0z_1\), and \(t=z_0z_1z_2\); the associated nearby-cycle
sheaves and triple-point signs; the normalization--conductor computation on
\(W\); hypercohomology ranks; the degree-three specialization-with-supports
and Gysin argument; and integral saturation.  Its output is the integral
isomorphism $\operatorname{sp}^q$ to $T_0$-invariants in every degree
\(0\le q\le4\), not an unproved substitution of the collapse pullback.
}

\ResultHead{Sections 7.5, 7.7 and 8: two load-bearing topology calculations}
{\ResultSmall
For Leray, the audit checks the cusp stalk, the indices at the fibres of
multiplicity three and four, local lift obstructions on three based
meridians, Mayer--Vietoris on the base, multiplicativity, and the exterior
pairings producing the two higher transgressions.  For van Kampen it checks
the explicit deck transformations, their third and fourth powers, the
zero-section base paths, the cusp disc, common relative signs, and the
presentation controlled by $p=-1$.  It then checks the Hurewicz--Whitehead
and smooth recognition conclusion.  Reviewer B classifies these three
stages as \textbf{VERIFIED--RECONSTRUCTED}.
}

\ResultHead{Section 9 and Section 10}
{\ResultSmall
The fibrewise period calculation, very-general N\'eron--Severi group,
signature $(1,1)$, exclusion of torus algebraic-reduction dimensions one
and two, and meromorphic descent are classified
\textbf{VERIFIED--INDEPENDENT}.  The Section~10 conductor mechanism is
classified \textbf{VERIFIED--LOCAL}: in the $xy=0$ and $xyz=0$ models,
the ambient differential has coefficients $(y,x,0)$ and $(yz,xz,xy)$,
vanishes on the conductor, and can map to conductor-supported intrinsic
torsion.  Its global deployment still uses topology, Picard information,
duality, and top-degree base change.
}

\ResultHead{Sections 3--6: no identified defect is not formal verification}
{\ResultSmall
Reviewer B read the period family, toric/logarithmic fillings, actions,
quotient atlases, compactness, Hausdorffness, and holomorphic gluing for
coherence; it reproduced the finite transformation and freeness arithmetic
and found no construction-level defect.  Its exact classification is
\textbf{REVIEWED---NO IDENTIFIED DEFECT}, because it did not formally
mechanize every analytic estimate and quotient-topology argument.
}

\ResultHead{The focused verification requests}
{\ResultSmall
\begin{enumerate}[leftmargin=1.35em,itemsep=1.5pt,topsep=1.4pt]
\item Check every proper-discontinuity, Hausdorffness, compactness, and
holomorphic-gluing hypothesis in Sections~3--6 lemma by lemma.
\item Audit the cited nearby-cycle, specialization-with-supports, Gysin,
duality, and base-change theorems with their exact hypotheses, naturality,
signs, and integral conclusions.
\item Complete or refute the explicit cusp-collapse map stratum by stratum
and derive its integral cellular matrices.  This remains relevant to the
first topology proof but is nonblocking for the nearby-cycle/Leray route.
\item Audit the global Picard and literature dependencies in Section~10
separately from its verified local conductor calculation.
\end{enumerate}
}

\vfill
{\fontsize{6.7pt}{7.7pt}\selectfont\color{AuditProvenance}
Evidence in this project is indexed in the peer-review map under
``Reviewer B --- 25 August corrective Pro-mode audit.''
The exact 87 MB evidence container has SHA-256
\texttt{155cf2aa4d243a92cb6a19d28a7ee3412c10de871b5283a4d7a96564fa735821};
its curated reports, ledgers, receipts, and executable checks are included.}
\endgroup

\newpage
\thispagestyle{empty}
\begingroup
\raggedright
\setlength{\parindent}{0pt}
\setlength{\parskip}{0pt}
\setlength{\fboxsep}{5pt}
\newcommand{\ResultHead}[1]{\par\vspace{3.4pt}{\fontsize{9.15pt}{10.25pt}\selectfont\bfseries\color{AuditReviewed}#1}\par\vspace{1.2pt}}
\newcommand{\ResultSmall}{\fontsize{8.0pt}{9.22pt}\selectfont}

{\fontsize{15.2pt}{16.6pt}\selectfont\bfseries Subsequent Codex validate/repair/disprove results}
\par\vspace{1.5pt}
{\fontsize{7.7pt}{8.9pt}\selectfont\color{AuditProvenance}
These are later workbench proof classifications, presented separately so
they are not attributed retrospectively to Reviewer B.  Each item below is
marked \textbf{VERIFIED} in the cited proof file as of 26 August 2026.}

\ResultHead{Global analytic threefold and all punctured-overlap maps --- VERIFIED}
{\ResultSmall
The workbench reconstructs the period family, finite logarithmic quotients,
toric cusp quotient, and three punctured-disc biholomorphisms.  In cusp
coordinates the covering map is
\[
 (z,\zeta_1,\zeta_2)\longmapsto
 (e^{2\pi i\zeta_1},e^{2\pi i\zeta_2},e^{2\pi i s}),
\]
with deck kernel \(\Lambda_{\rm tor}\times\langle g_0\rangle\) and residual
deck group \(\bar\Lambda\).  The quotient proof establishes a connected
compact Hausdorff complex threefold and a proper surjective holomorphic map
\(f:X\to\mathbb P^1\) with the stated cusp and multiple fibres.  See
\texttt{research/candidate\_geometry.tex}.
}

\ResultHead{Intrinsic specialization and the Leray differentials --- VERIFIED}
{\ResultSmall
The workbench gives an integral, stalkwise proof of the same
\(\operatorname{sp}^q\)-isomorphisms, including the supported
specialization/Gysin comparison and finite-fibre stalk maps.  It then closes
the formal identification of the local lift cocycle with Leray $d_2$: for
\(K=Rf_*\mathbb Z_X\), the canonical truncation triangle is
\[
 R^0f_*\mathbb Z_X\longrightarrow\tau_{\le1}K
 \longrightarrow R^1f_*\mathbb Z_X[-1]
 \xrightarrow{\kappa}R^0f_*\mathbb Z_X[1],
\]
and $H^0(\kappa):H^0(R^1f_*\mathbb Z)\to H^2(R^0f_*\mathbb Z)$ is the
Leray transgression.  Functorial Mayer--Vietoris for this triangle computes
it by the three annular lift differences.  With the workbench orientation,
\[
 d_2(12\gamma)=\omega,\qquad
 d_2(2q)=\omega[\delta],\qquad
 d_2(2\gamma uw)=\omega[\gamma\delta],
\]
so $H^1=H^2=H^3=0$ and $H_2=0$.  See
\texttt{research/integral\_leray.tex}.
}

\ResultHead{Based van Kampen and sphere recognition --- VERIFIED}
{\ResultSmall
Using one zero-section basepoint and fixed clockwise meridians, the workbench
derives the actual gluing presentation
\[
 \pi_1(X)=\langle c,x,y\mid c\ {\rm central},\ xy=1,\ x^3=c,\ y^4=c^{-1}\rangle=1.
\]
It proves the cusp meridian bounds the toric holomorphic disc and that no
unrecorded fibre drift appears.  Universal coefficients, Poincar\'e duality,
Hurewicz, Whitehead, and smooth six-sphere recognition then give
\[
 X\cong_{\rm diff}S^6,
 \qquad \langle c_3(T_X),[X]\rangle=\chi(X)=2.
\]
See \texttt{research/integral\_leray.tex}.
}

\ResultHead{Algebraic reduction --- VERIFIED}
{\ResultSmall
The candidate's period argument is expanded by an explicit Appell--Humbert
matrix bridge and a complete meromorphic-descent proof.  The result is
\[
 \mathcal M(X)=f^*\mathbb C(t),\qquad a(X)=1,
\]
with $f$ the algebraic reduction.  This proof is independent of the
topology and CDP comparison.  See
\texttt{research/algebraic\_dimension.tex}.
}

\ResultHead{Proposition 10.8: uniform non-torsion --- VERIFIED}
{\ResultSmall
The workbench now gives a complete proof of the candidate's
Proposition~10.8 branch:
\[
 f_*\mathcal O_X=\mathcal O_B,\qquad
 R^1f_*\mathcal O_X\simeq\mathcal O_B\oplus\mathcal O_B(-1),\qquad
 R^2f_*\mathcal O_X\simeq\mathcal O_B(-1),\qquad
 \operatorname{Pic}(X)\simeq\mathbb C.
\]
All base-change maps in degrees \(q=0,1,2\) are isomorphisms at every
scheme-theoretic fibre; moreover
\(\operatorname{ord}(\mathcal O_X(S_j)|_{S_j})=m_j\), with
\((m_1,m_2)=(3,4)\).  For
\(\lambda=a\gamma+bu+cw+d\delta\), the exact condition that its Dolbeault
image lie in the global Picard line is
\[
 c\tau+d\mu-b=0.
\]
The independence of \(1,\tau,\mu\) leaves only \(\mathbb R\) and a countable
exceptional set, so every remaining parameter gives a line bundle
non-torsion on every full fibre and on \(S_1,S_2\).  The cusp extension uses
intrinsic nearby cycles, not the still-\textbf{OPEN} optional collapse comparison.
This closes Proposition~10.8 within the reconstructed construction; the
cellular-collapse/Mayer--Vietoris comparison and external CDP adjudication
remain \textbf{OPEN}.  See \texttt{research/analytic\_picard.tex}.
}

\ResultHead{Provenance and the remaining open route}
{\ResultSmall
The construction, nearby-cycle strategy, transgression arithmetic,
fundamental-group presentation, and period strategy originate in the
candidate manuscript.  Codex supplied proof-level reconstructions and, in
particular, the explicit cusp-overlap quotient, the canonical-truncation
identification of the \v{C}ech obstruction with Leray $d_2$, a full based
coordinate van Kampen proof, the Appell--Humbert bridge, and the omitted
meromorphic-descent details.  The separate explicit cusp-collapse map,
its proposed equality with $\operatorname{sp}^q$, and the first cellular
Mayer--Vietoris matrices remain \textbf{OPEN}.  They are geometric claims
worth resolving, but they are not premises of the verified non-collapse
recognition route.
}

\vfill
{\fontsize{6.65pt}{7.65pt}\selectfont\color{AuditProvenance}
The earlier review and page-by-page annotations follow as a preserved
historical record.  Where an older status label conflicts with this front
matter, consult the current proof files and ledgers in the separately
maintained public workbench; the change is recorded rather than silently
retrofitted into an earlier reviewer's voice.}
\endgroup
